\section{Steenrod operations on homology}

Let $C$ be a chain complex, over $R = \mathbb{F}_p$. Consider $C$ as a $R\cyc_p$-complex with trivial $\cyc_p$-action and $C^{\otimes p}$ with the permutation action. Let $H(C)^{(p)}$ be the `inflation'' of the graded module $H(C)$:
\[
H(C)^{(p)}_k = \begin{cases}
	H(C)_j & \text{ if $k = pj$} \\
	0 & \text{ otherwise.}
\end{cases}
\]
\begin{lemma}
	There are isomorphisms \martin{I rewrote the notation for the Lemma, since Tate homology is for groups, not spaces or chain complexes.}
	\begin{align*}
		\widehat{H}(\cyc_p;C)&\cong H(C)\ot \widehat{H}(\cyc_p;\bF_p)
		\\
		\widehat{H}(\cyc_p;C^{\ot p}) & \cong H(C)^{(p)}\ot \widehat{H}(\cyc_p;\bF_p) 
	\end{align*}
\end{lemma}
\martin{For the second isomorphism the proof right now only gives a spectral sequence from the term on the right to the term on the left (see next comments).}
\begin{proof}
	The first isomorphism is standard, because $\cyc_p$ acts trivially on $C$. For the second, recall that the bicomplex $C\ot_{\cyc_p} \widehat{W}(p)$ gives rise to two spectral sequences. In one of them, the first page is obtained by taking homology only with respect to the differential of $C$, and it is isomorphic to $H(C)^{\ot p}\ot_{\cyc_p} \widehat{W}(p)$ by the Künneth isomorphism.

	Now, if $M$ is a (graded) $R\cyc_p$-module, we can decompose it as $M = M'\oplus M''$, where $M'$ has a trivial action and $M''$ has a free action (it is an induced module), and there is an isomorphism $H(M\ot_{\cyc_p} \widehat{W}(p);\bF_p)\cong M^\prime\ot\widehat{H}(\cyc_p;\bF_p)$. \martin{Instead of using a spectral sequence we could just use an argument similar to this one from the beginning. Note that this is quite similar to Lemma 6.}
Since the part with trivial action in $H(C)^{\ot p}$ is isomorphic to $H(C)^{(p)}$ (if we choose a basis of $H(C)$, then we can identify these two with the map $x\mapsto x^{\ot p}$), we conclude that the second page is isomorphic to $H(C)^{(p)}\ot \widehat{H}(\cyc_p;\bF_p)$. \martin{It remains to show that the spectral sequence converges at the second page. I don't think that this is immediate, since the pages $pk+1$ may have non trivial differentials.}
\end{proof}

If $X$ is a simplicial set, a natural equivariant map
\[
\chains(X)\longrightarrow \chains(X)^{\ot p}\ot \widehat{W}(p)
\]
induces, after taking Tate homology, a map
\[
H(X)\longrightarrow H(X)^{(p)}\ot \bF_p[\{e_i\}]
\]
Thus, decomposing the target into summands indexed by $i$:
\[
x\mapsto \sum_i f_i(x)\ot e_i
\]
defines operations $x\mapsto f_i(x)$. I think that $f_i$ is the $D_i$-operation (so $D_{k(p-1)}(x)$ computes $P^0(x)$).